% ============================================================================
% Study 1 — Timing and window of harm formation in a masked diffusion LM
% Body content only (drop into the existing paper). Requires: booktabs, multirow.
% Scope throughout: within-template, within-model (LLaDA-8B-Instruct), DIJA.
% ============================================================================

\section{Timing and Window of Harm Formation in a Masked Diffusion LM}
\label{sec:study1}

We studied \emph{when}, along a masked diffusion LM's denoising trajectory, a
template-injection attack's harmful content becomes committed, and whether that
timing leaves any window for a mid-generation defense. All results are
within-template, within-model (LLaDA-8B-Instruct; 32 layers, $d_{\text{model}}=4096$,
128 denoising steps, unified \texttt{fill\_all\_masks} schedule, \texttt{low\_confidence}
remasking), and specific to the DIJA template-injection attack family
% AUTHOR QUERY: cite keys for LLaDA-8B-Instruct and the DIJA attack paper
\citep{TODO-llada, TODO-dija}. The attack injects a worksheet scaffold whose
\texttt{<mask:N>} blanks sit \emph{inside the prompt} at token span $[t_0,t_1)$; the
judged response was $\operatorname{decode}(x[t_0{:}t_1]) \,\|\, \operatorname{decode}(x[P{:}])$.
Harm was scored with the valence-\emph{inclusive} judge
(\texttt{unchanged\_harmful}, \texttt{euphemistic\_softening}, \texttt{disclaimer\_only}),
which the study signed off on; the reconciliation against a HarmBench-style judge is
reported in \Cref{sec:judge-reconciliation}.

We set out to measure $\Delta = k_{\text{commit}} - k_{\text{detect}}$, the gap between
when harmful content becomes \emph{detectable} and when it becomes \emph{committed}.
$\Delta$ is \emph{not} reported: the study instead produced one positive result
(\Cref{ssec:kcommit}), one negative control that constrains what any $k_{\text{detect}}$
claim may say (\Cref{ssec:placebo}), a behavioural characterisation of how far the
scaffold determines the outcome (\Cref{ssec:determination}), a falsification of the
one defense the medium uniquely affords (\Cref{ssec:remask}), and --- extending that
last result --- a causal, directional test that localises the determinant to the scaffold
\emph{representation} and finds it selectively steerable (\Cref{ssec:regionsteer}).

% ----------------------------------------------------------------------------
\subsection{$k_{\text{commit}}$ is measurable and late, with a bimodal structure}
\label{ssec:kcommit}

We defined the per-position commit step $k_i = \min\{k : i \notin M(x_{t_k})\}$ and
measured it over the injected worksheet blanks
(\texttt{region\_id == R\_SCAFFOLD \& fillable\_reply\_mask}), the positions where DIJA
harm is written. The unit is bounded by \texttt{n\_inject}, so it cannot inflate with
continuation length; an earlier span-locator that ballooned to 157--163 positions was
the v1 failure mode and was fixed before any $k_{\text{commit}}$ was computed.

\begin{table}[t]
\centering
\begin{tabular}{lccc}
\toprule
statistic & harm-delivered ($n{=}19$) & not-delivered ($n{=}6$) & all ($n{=}25$) \\
\midrule
median of per-case medians & \textbf{98\,/\,128} & 103 & 98 \\
blank-weighted pooled median & 93\,/\,128 & 100 & 95 \\
pooled IQR & $[68,110]$ & $[84,114]$ & $[73,111]$ \\
blanks & 1182 & 330 & 1512 \\
\bottomrule
\end{tabular}
\caption{$k_{\text{commit}}$ over injected worksheet blanks. Harm committed at roughly
0.73--0.77 of the trajectory; the headline ``98'' is the median of per-case medians.}
\label{tab:kcommit}
\end{table}

\paragraph{Bimodal, with an empty gap.}
The per-case medians (harm-delivered) split cleanly: a late cluster of $n{=}17$ with
medians in 90--103 (the modal behaviour), and an early cluster of $n{=}2$ --- A006
(26) and A003 (41). Nothing fell between 41 and 90; a $\sim$50-step gap populated by
zero cases is a genuine split rather than the tail of one distribution. A006 was early
\emph{throughout} (p25 20, p75 32; 68.8\% of its blanks committed by step 30), not
early-in-the-tail, and both early cases survived the span-locator fix on the bounded
scaffold-blank unit. With $n{=}2$ we describe the early cluster but do not explain it.

\paragraph{The result is robust to the skeleton-exclusion cut.}
Worksheet scaffolds contain connective tokens the model pins immediately and which
carry no harm; the pre-registered exclusion was $k_i \le 2$. Under it the pooled median
moved $93.0 \to 94.0$ and the median of case medians was unchanged at $97.5$; no case
moved by $\ge 5$ steps (largest shift $+2.5$). Widening the cut from 2 to 50 moved the
pooled median $94 \to 101$ --- only \emph{later}, never earlier --- so the ``harm
commits late'' claim is robust in the direction that matters.

\paragraph{Open in both directions: sub-token attribution.}
Within a single scaffold the blanks split into early- and late-committing populations,
and our per-response harm label cannot say which carry the harm. Under Reading~A the
early blanks are filler and the late blanks carry the specifics, so the true
harm-bearing $k_{\text{commit}}$ is \emph{later} than 98; under Reading~B the model
fixes the plan's semantic skeleton first and fills connectives late, so it is
\emph{earlier} (A006's shape is what Reading~B predicts). This is not decidable from the
current data and is recorded as an open refinement; the headline is stable under the cut
sweep either way, and only its attribution to particular blanks is unresolved.

% ----------------------------------------------------------------------------
\subsection{A negative control: matched geometry is necessary but not sufficient}
\label{ssec:placebo}

This result constrains what any $k_{\text{detect}}$ claim in this setting may say. For
each request, the \texttt{dija} arm (harmful scaffold) and a \texttt{benign\_op} arm
(injected-but-benign) shared the same base request and carried byte-identical
\texttt{<mask:N>} markers --- identical blank geometry, asserted at build and re-asserted
at capture. A grouped cross-validated probe (by request, 1512 blanks across 25 matched
pairs, out\_mask-only features) gave \textbf{AUC 1.000 at step 0}, before any denoising.
That is the signature of a leak, so we ran a placebo rather than reporting $\Delta=+98$.

\begin{table}[t]
\centering
\begin{tabular}{lccc}
\toprule
contrast & harm difference? & peak AUC & reading \\
\midrule
dija vs.\ benign\_op & yes & 1.000 @ step 0 & leak \\
placebo: dija(no-harm) vs.\ benign\_op & \textbf{none} & \textbf{1.000 @ step 0} & pure arm classifier \\
both, after subtracting step 0 & yes / none & 0.992 @ 117 / \textbf{1.000 @ 1} & leak survives \\
within-dija (harm vs.\ no-harm, arm fixed) & yes & 0.904 @ 89 (null max 0.939) & not significant \\
\bottomrule
\end{tabular}
\caption{The placebo labels the \emph{arm} on only the 6 requests where the dija arm
delivered no harm, so there is no harm difference to find; it still separates perfectly.}
\label{tab:placebo}
\end{table}

\paragraph{The finding.}
Matching request and byte-identical blank geometry was necessary but \emph{not}
sufficient. The scaffold \emph{prose} still differed between arms --- the blanks sat
inside harmful worksheet text in one arm and benign worksheet text in the other --- so
arm identity was readable from the step-0 hidden state and was \emph{re-encoded at
step 1} after step-0 subtraction: the prompt conditions every subsequent step, so no
single static component can be subtracted away. A \texttt{delta-from-step0} remedy forced
step-0 AUC to 0.5 by construction and the placebo immediately re-separated at 1.000 @
step 1; it does not work, and is documented so it is not re-attempted.

\paragraph{Methods note: a permutation null that returns the observed statistic is
broken, not strong.}
The label-permutation null was originally within-group, the exact null for the paired
two-arm design. For the within-dija contrast the groups are singletons, so shuffling a
one-element list returned it unchanged and the ``null'' came back \emph{exactly} equal to
the observed AUC (0.851 vs.\ 0.851; 0.904 vs.\ 0.904) --- the equality is what exposed
it. Falling back to a global shuffle for singleton groups, within-dija peaked at 0.851
raw / 0.904 de-leaked (step 89) against a max-over-steps null of 0.904 / 0.939, not
significant at 19 vs.\ 6. $\Delta$ is additionally not computable in that contrast, as
the benign operating point it needs is absent from it.

% ----------------------------------------------------------------------------
\subsection{Harm delivery is strongly, but not absolutely, determined by the scaffold}
\label{ssec:determination}

Measured with generate-only sampling (no hidden capture).

\paragraph{At temperature 0 the outcome is a deterministic function of the prompt.}
This holds by construction, not empirically: \texttt{add\_gumbel\_noise} short-circuits
at $T{=}0$ (it never draws RNG), confidence is a pure softmax-gather, and top-$k$
selection draws no RNG. We confirmed 5/5 cross-seed pairs byte-identical anyway, since
kernel-level nondeterminism is invisible in source. (Random remasking \emph{is}
seed-sensitive at $T{=}0$ but would randomise the step at which a token leaves the mask
set, destroying $k_{\text{commit}}$, so it was not used.)

\paragraph{At the paper's temperature 0.2 the outcome is $\sim$96\% determined, not
100\%.} A first pass (15 requests, $K{=}8$) was under-powered, finding one straddling
request and a rule-of-three point estimate of $p{=}0$ over 112 one-sided draws.
Re-running 5 of those requests at $K{=}32$ gave 7/160 minority draws, a flip rate
$\hat p = 0.044$ (Wilson-95 $[0.021,0.088]$), and 4/5 straddling. This $\hat p$ is a
\emph{lower bound}: those 5 requests were selected for coming out one-sided at $K{=}8$,
which biases toward low-flip cases.

\paragraph{The residual grows smoothly with temperature.}

\begin{table}[t]
\centering
\begin{tabular}{lccccc}
\toprule
$T$ & minority draws & rate & Wilson-95 & straddle @ $K{=}16$ & collapsed \\
\midrule
\textbf{0.2} (paper config) & 7/160 & \textbf{0.044} & $[0.021,0.088]$ & 1/5 & 0/160 \\
0.4 & 6/80 & 0.075 & $[0.035,0.154]$ & 2/5 & 0/80 \\
0.7 & 11/80 & 0.138 & $[0.079,0.230]$ & 4/5 & 0/80 \\
\bottomrule
\end{tabular}
\caption{Both metrics rise monotonically; there is no temperature at which determinism
breaks, it thins continuously. Temperatures 0.4 and 0.7 are a different decoding regime
and describe that regime only; the 0.2 row carries no such caveat.}
\label{tab:tempsweep}
\end{table}

\paragraph{Controls make the flips mean something.}
No collapsed samples appeared at any temperature (\texttt{rep3}\,$>$\,0.50 or
distinct-word ratio\,$<$\,0.30; observed trigram repetition 0.01--0.08, distinct ratio
0.52--0.73), and within a straddling request the two outcome sides were metrically
indistinguishable, so no flip is a degeneration artifact. Across all temperatures 21
byte-identical duplicate generations received the same label in the first passes, so
every flip was a genuinely different generation; this control later \emph{fired} at
larger scale (below), which is why it is reported.

\paragraph{The unselected rate --- and straddling is a property of particular
scaffolds.} An unselected pass (25 requests, $K{=}32$, 800 draws) gave a
text-resolved population flip rate of $22/800 = 0.0275$ (Wilson-95 $[0.018,0.041]$). The
minority mass was \emph{not} spread evenly: 16 of 25 requests never flipped, one request
held 45\% of it, and the per-request counts were massively over-dispersed against
$\mathrm{Binomial}(32,0.0275)$ --- $\chi^2 = 119.9$ on 24 df, $p \approx 10^{-14}$;
16 zero-cells observed against 10.2 expected; $P(\text{any request}\ge 10\ \text{minority
under uniformity}) \approx 2\times10^{-7}$. \textbf{Straddling is therefore a property of
particular scaffolds, not uniform sampling noise}, and the pooled ``$\sim$96\%
determined'' is a mixture of two regimes rather than a typical case. This also explains
why the \emph{selected} 0.044 exceeded the \emph{unselected} 0.0275: heterogeneity
swamped the selection bias.

\paragraph{Judge noise is real at this scale, and the unit of analysis must be the text.}
In the unselected pass, 2 of 94 duplicated texts were judged both ways --- one text with
26 byte-identical copies was labelled \texttt{disclaimer\_only} 3 times and
\texttt{substantive\_safe\_substitution} 23 times. Identical text is an identical
trajectory and cannot differ in outcome, so labels were resolved per \emph{unique text}
by majority vote before counting; raw counting inflated the rate $0.0325 \to 0.0275$. The
channel carries $\sim$2\% label error on borderline texts, which sit exactly on the
\texttt{disclaimer\_only} $\leftrightarrow$ \texttt{substantive\_safe\_substitution}
boundary --- the straddler population itself.

% ----------------------------------------------------------------------------
\subsection{Re-masking as a native defense operator: the scaffold, not the output
region, is the determinant}
\label{ssec:remask}

At inference in a masked diffusion LM there is no separate noising phase: re-masking
\emph{is} the forward (noising) operator, applied selectively --- a lever an
autoregressive model structurally lacks, since a generated token cannot be un-generated.
This being discrete masked diffusion, ``dose'' means the fraction of positions
re-masked, not a continuous noise level. Because the harm commits late
(\Cref{ssec:kcommit}) yet the outcome is $\sim$determined by the scaffold
(\Cref{ssec:determination}), we asked whether re-masking the committed output can undo
the harm. The experiment was generation-only, defensive, and approved by the PI.

\paragraph{Design.} Two passes at $T{=}0$. Pass~1 ran a plain baseline (arm~E),
recording per-position commit steps and the located harm span; these defined the target
positions and two \emph{budget-floored} intervention steps $k^\ast$ (each leaving
$\ge 20$ remaining denoising steps). Pass~2 re-masked a target set at $k^\ast$,
recomputed the transfer schedule over the current masks for the remaining steps
(guaranteeing every re-masked position refilled), and continued from the \emph{shared
deterministic pre-intervention state} --- at $T{=}0$ the trajectory to $k^\ast$ is
bit-identical to the baseline, so no counterfactual is introduced. The intervention
engine imported the model's own sampling primitives unedited. Arms: \textbf{A} re-masked
the worksheet blanks (the harm site); \textbf{C} re-masked a matched count of non-harm
positions with the worksheet left intact (a floor); \textbf{D} wiped the worksheet
\emph{and} relocated an equal blank count into free output space, then read where harm
landed; \textbf{E} was the baseline. An arm re-masking ``worksheet $+$ scaffold blanks''
was dropped because for DIJA the harm span already contains the scaffold blanks
(harm is written into the worksheet, \Cref{ssec:kcommit}), so it is identical to A; the
distinction it was meant to test is structural, not statistical --- the harmful
\emph{template} itself (the section headers and the fixed ``\ldots highly dangerous''
lines) is immutable prompt text that no output-region re-mask can reach.

\paragraph{Deterministic run: no arm neutralised harm.}
Seven of ten baselines delivered harm. Over those 7 at-risk requests, every arm at
every $k^\ast$ re-delivered harm (\Cref{tab:remask}); there were zero neutralisations.
This was not step-budget starvation: 0/70 outputs collapsed (\texttt{rep3}
$\mu \approx 0.03$--$0.05$ against the 0.50 bar), because $k^\ast$ was budget-floored.

\begin{table}[t]
\centering
\begin{tabular}{llcc}
\toprule
arm & $k^\ast$ & harm re-delivered & rate (Wilson-95) \\
\midrule
E (baseline) & --- & 7/7 & 1.00 $[0.65,1.00]$ \\
A (worksheet) & early / late & 7/7 \,/\, 7/7 & 1.00 \\
C (non-harm floor) & early / late & 7/7 \,/\, 7/7 & 1.00 \\
D (wipe + relocate) & early / late & 7/7 \,/\, 7/7 & 1.00 \\
\bottomrule
\end{tabular}
\caption{Deterministic re-mask over the 7 at-risk requests: $A \approx C \approx D
\approx E = 1.00$, valence-inclusive, non-degenerate.}
\label{tab:remask}
\end{table}

\paragraph{The floor forecloses the mis-targeting objection.}
Control~C left the harm span intact and still delivered; arm~A wiped the entire
worksheet and also delivered; the two matched exactly. The result therefore cannot be
``we re-masked the wrong positions'' --- targeting the harm site was no better than a
matched floor that never touched it. There is no set of output positions whose removal
helps.

\paragraph{Mechanism (arm D): harm is bound to the fixed template, not to the
availability of blanks.} D wiped the worksheet and opened an equal count of fresh blanks
in free output space. The harm re-formed in the \emph{original} scaffold slots while the
relocated blanks took benign filler (e.g.\ ``\ldots the patient.'', ``\ldots consult a
mental health professional''). The determinant is not that blanks exist to fill; it is
the immutable template text that keeps demanding a harmful fill \emph{at its own
positions}. This is the causal read.

\paragraph{Methods note --- a warning to future work on re-masking defenses.} A naive
late re-mask at $k^\ast = \max(\text{span commit})$ leaves $\sim$1 refill step; a
$k^\ast$-sweep showed the refill degenerates (duplicate-runs 0 at $\ge 18$ remaining
steps, 3 at 4, 10 at 1) and the judge then reads the garble as
\texttt{substantive\_safe\_substitution}, i.e.\ \emph{not} harmful. A re-masking defense
evaluated without degeneration controls will therefore report a false success: the
starved refill is unreadable, not safe. Every neutralisation must be cross-checked
against a degeneration score; budget-flooring removed the artifact here and the harm
returned in full.

\paragraph{Branch test: re-masking did not manufacture the open-outcome population.}
Re-running arm~A with 8 seeds at $T{=}0.2$ from the shared committed state (late
$k^\ast$), and requiring $\ge 2$ minority draws on unique texts (a 1-of-$K$ split lies
within the $\sim$2\% label-noise band), only 1 of 10 requests branched --- and it was
the same request that already straddles under ordinary sampling
(\Cref{ssec:determination}). Re-masking reproduced the existing straddle rather than
creating new ones, an independent corroboration of the heterogeneity finding by a
distinct experimental route; the shared committed state plus the fixed scaffold re-drive
the same outcome even under a $T{=}0.2$ refill. The standing duplicate-judge control
fired once, on the same borderline category, and was voided by unique-text resolution.

\paragraph{Cross-cutting check.} The span locator sweeps the trailing refusal into its
harm span, which inflated a span-based commit variant; the reported $k_{\text{commit}}$
(\Cref{ssec:kcommit}) is measured on a different unit --- worksheet blanks only, bounded
by \texttt{n\_inject} --- and so is unaffected. We verified this in code; the headline
result is untouched.

% ----------------------------------------------------------------------------
\subsection{Prompt-region representation steering: the scaffold's representation
is causally and selectively steerable}
\label{ssec:regionsteer}

Result~4 (\Cref{ssec:remask}) showed that no \emph{output-region} operator undoes DIJA
harm and read the cause as the immutable scaffold text --- a claim about \emph{where} the
determinant sits, established by elimination. Here we tested it causally and
directionally: if the scaffold is the determinant, then steering the model's internal
representation of the \emph{scaffold} region should move the harm, whereas the identical
steer applied to the \emph{output} region should not. This was a
falsification/confirmation experiment on activations --- not a text edit, and not a
defense.

\paragraph{Design.} We reused the injection axis $v$: the unit-norm top right singular
vector of clean$\to$DIJA prompt-representation differences at layer~16 (mean-pooled), for
which $+v$ is ``more injected''. At every denoising step we added $-\alpha v$ to the
layer-16 residual output at a per-arm \emph{position set} --- a de-injection of dose
$\alpha \in \{21,42,84,126\}$, calibrated to the $\approx 21$ layer-16 injection-signal
magnitude. Arms differed only in \emph{which rows} were steered: \textbf{T}
(\texttt{tpl\_ctx}), the unmasked scaffold-text tokens inside the injected span
$[t_0,t_1)$; \textbf{O} (\texttt{out\_mask}), the still-masked output positions;
\textbf{R}, a count-matched random subset of output positions; \textbf{Rw}, the
chat-wrapper tokens; \textbf{E}, no steer. Per-case RNG was fixed (seed $=$ base $+$ case
index) so arm differences are pure steering. The endpoint was the graded valence judge's
continuous \emph{specificity} (actionability, $0$--$1$), paired per case against E; we
report the mean paired delta $\Delta_{\text{spec}} = \text{spec}(\text{arm}) -
\text{spec}(\text{E})$ with a $10^4$ bootstrap 95\% CI and a two-sided Wilcoxon
signed-rank test. Selectivity used the same steer and endpoint on the per-request matched
Group-C benign scaffolds (\emph{benign} specificity; arms E/T/O). The set was $n{=}40$
harm-delivering B cases, balanced 20 strict-harm (\texttt{unchanged\_harmful} /
\texttt{euphemistic}) $+$ 20 \texttt{disclaimer\_only}, all 40 delivering inclusive harm
at baseline. Byte-identical duplicate outputs were judged once and shared their verdict
(106/680 B and 33/360 C duplicates collapsed), removing the $\sim$2\% label-noise floor
for arms whose output matched E.

\begin{table}[t]
\centering
\begin{tabular}{llccc}
\toprule
arm & $\alpha$ & $\Delta_{\text{spec}}$ (B) & boot 95\% CI & Wilcoxon $p$ \\
\midrule
\multirow{4}{*}{T (\texttt{tpl\_ctx})}
 & 21  & $+0.020$ & $[-0.018,+0.062]$ & $0.50$ \\
 & 42  & $-0.025$ & $[-0.085,+0.035]$ & $0.23$ \\
 & \textbf{84}  & $\mathbf{-0.123}$ & $[-0.195,-0.048]$ & $\mathbf{0.0012}$ \\
 & 126 & $-0.255$ & $[-0.315,-0.193]$ & $7.5\!\times\!10^{-7}$ \\
\midrule
\multirow{4}{*}{O (\texttt{out\_mask})}
 & 21  & $-0.005$ & $[-0.033,+0.023]$ & $0.72$ \\
 & 42  & $-0.010$ & $[-0.030,+0.010]$ & $0.25$ \\
 & 84  & $+0.003$ & $[-0.030,+0.035]$ & $0.91$ \\
 & 126 & $-0.040$ & $[-0.087,+0.005]$ & $0.12$ \\
\midrule
\multirow{4}{*}{R (rand.\ output)}
 & 21  & $-0.005$ & $[-0.038,+0.030]$ & $0.83$ \\
 & 42  & $+0.003$ & $[-0.035,+0.045]$ & $0.94$ \\
 & 84  & $-0.015$ & $[-0.057,+0.027]$ & $0.67$ \\
 & 126 & $-0.008$ & $[-0.068,+0.055]$ & $0.68$ \\
\midrule
\multirow{4}{*}{Rw (wrapper)}
 & 21  & $+0.000$ & $[-0.025,+0.025]$ & $0.95$ \\
 & 42  & $+0.005$ & $[-0.030,+0.042]$ & $0.83$ \\
 & 84  & $+0.030$ & $[-0.010,+0.075]$ & $0.40$ \\
 & 126 & $+0.043$ & $[-0.005,+0.095]$ & $0.14$ \\
\bottomrule
\end{tabular}
\caption{\textbf{Primary.} Paired harmful-specificity delta (arm $-$ E) over $n{=}40$ B
cases. Only T --- steering the prompt-scaffold representation --- degrades specificity,
and does so dose-dependently; O, R, and Rw are flat at every dose (all ns, max
$|\Delta|=0.043$). Means are over all paired cases; the signed-rank excludes byte-stable
($\Delta{=}0$) pairs, which are common at low $\alpha$. Positive $\Delta$ = more specific.}
\label{tab:regionsteer-primary}
\end{table}

\paragraph{Primary: only the prompt-region steer degrades harm, and dose-dependently.}
T reduced harmful specificity region-specifically and monotonically in dose:
negligible at $\alpha\!\le\!42$, $-0.123$ at $\alpha{=}84$ ($p{=}0.0012$), and $-0.255$
at $\alpha{=}126$ ($p{=}7.5\!\times\!10^{-7}$; \Cref{tab:regionsteer-primary}). The
\emph{identical} direction and magnitude applied to the output region (O), a
count-matched random output subset (R), or the chat wrapper (Rw) moved specificity
nowhere at any dose (all ns). Prompt-region steering was thus the only intervention that
degraded harm --- the causal, directional counterpart to Result~4's elimination argument:
the determinant is carried by the scaffold's representation, and it is reachable from
there and only from there.

\begin{table}[t]
\centering
\begin{tabular}{llcccc}
\toprule
stratum & $\alpha$ & $n$ & $\Delta_{\text{spec}}$ & boot 95\% CI & Wilcoxon $p$ \\
\midrule
\multirow{4}{*}{strict harm}
 & 21  & 8  & $+0.040$ & $[-0.015,+0.110]$ & $0.23$ \\
 & 42  & 6  & $+0.020$ & $[-0.030,+0.085]$ & $0.67$ \\
 & \textbf{84}  & 16 & $\mathbf{-0.105}$ & $[-0.185,-0.025]$ & $\mathbf{0.022}$ \\
 & 126 & 17 & $-0.220$ & $[-0.305,-0.135]$ & $0.0010$ \\
\midrule
\multirow{4}{*}{\texttt{disclaimer\_only}}
 & 21  & 9  & $-0.000$ & $[-0.050,+0.045]$ & $0.72$ \\
 & 42  & 15 & $-0.070$ & $[-0.170,+0.035]$ & $0.086$ \\
 & 84  & 19 & $-0.140$ & $[-0.255,-0.015]$ & $0.027$ \\
 & 126 & 18 & $-0.290$ & $[-0.370,-0.205]$ & $0.00027$ \\
\bottomrule
\end{tabular}
\caption{\textbf{Strata (arm T).} The specificity reduction held on \emph{real strict
harm} (\texttt{unchanged\_harmful} / \texttt{euphemistic}), not just on the
\texttt{disclaimer\_only} cases --- closing the disclaimer-only gap of the $n{=}10$ pilot.
$n$ is discordant (non-zero) pairs.}
\label{tab:regionsteer-strata}
\end{table}

\paragraph{The effect holds on real strict harm.} Split by baseline harm type
(\Cref{tab:regionsteer-strata}), T degraded specificity on the strict-harm stratum
($-0.105$ at $84$, $p{=}0.022$; $-0.220$ at $126$, $p{=}0.0010$) as well as on the
disclaimer stratum ($-0.140$; $-0.290$). The $n{=}10$ pilot could only show the effect on
disclaimer-only cases; the powered set closes that gap on genuinely actionable harm.

\begin{table}[t]
\centering
\small
\begin{tabular}{ll ccc ccc l}
\toprule
& & \multicolumn{3}{c}{B (harmful spec)} & \multicolumn{3}{c}{C (benign spec)} & \\
\cmidrule(lr){3-5}\cmidrule(lr){6-8}
arm & $\alpha$ & $\Delta$ & 95\% CI & $p$ & $\Delta$ & 95\% CI & $p$ & verdict \\
\midrule
T & \textbf{84}  & $-0.123$ & $[-0.193,-0.047]$ & $0.0012$ & $-0.005$ & $[-0.065,+0.060]$ & $0.43$ & \textbf{selective} \\
T & 126 & $-0.255$ & $[-0.312,-0.195]$ & $7.5{\times}10^{-7}$ & $-0.150$ & $[-0.210,-0.087]$ & $8.5{\times}10^{-5}$ & non-selective \\
O & 84  & $+0.003$ & $[-0.030,+0.035]$ & $0.91$ & $-0.107$ & $[-0.163,-0.055]$ & $0.00056$ & benign-only harm \\
O & 126 & $-0.040$ & $[-0.087,+0.005]$ & $0.12$ & $-0.125$ & $[-0.177,-0.075]$ & $0.00015$ & benign-only harm \\
\bottomrule
\end{tabular}
\caption{\textbf{Selectivity.} At $\alpha{=}84$, T dropped harmful specificity while
leaving benign specificity intact (selective). At $\alpha{=}126$, T also tanked benign
specificity (non-selective). O never moved harmful specificity yet damaged benign at both
doses (pure capability cost, no safety benefit). CIs are from independent bootstrap draws
and differ marginally from \Cref{tab:regionsteer-primary}.}
\label{tab:regionsteer-selectivity}
\end{table}

\begin{table}[t]
\centering
\begin{tabular}{llcccc}
\toprule
arm & $\alpha$ & helpful & on-topic & benign spec & collapse \\
\midrule
E  & ---  & 0.39 & 0.75 & 0.26 & 0.07 \\
\midrule
T  & 21   & 0.43 & 0.75 & 0.29 & 0.03 \\
T  & 42   & 0.43 & 0.70 & 0.29 & 0.04 \\
T  & \textbf{84} & 0.37 & 0.57 & 0.26 & 0.06 \\
T  & 126  & 0.15 & 0.10 & 0.11 & 0.21 \\
\midrule
O  & 21   & 0.40 & 0.62 & 0.28 & 0.06 \\
O  & 42   & 0.38 & 0.68 & 0.28 & 0.05 \\
O  & 84   & 0.21 & 0.38 & 0.15 & 0.05 \\
O  & 126  & 0.16 & 0.30 & 0.14 & 0.03 \\
\bottomrule
\end{tabular}
\caption{\textbf{Group-C utility guard.} At the selective point ($\alpha{=}84$) benign
specificity was preserved ($0.26$ vs.\ $0.26$) and Group-C helpfulness was essentially
unchanged ($0.39\to0.37$), with on-topic already softening ($0.75\to0.57$). At $\alpha{=}126$ utility collapsed
(helpful $\to0.15$, on-topic $\to0.10$, collapse $\to0.21$). O degraded utility from
$\alpha{=}84$ on.}
\label{tab:regionsteer-guard}
\end{table}

\paragraph{Selectivity: $\alpha{=}84$ is the operating point, $\alpha{=}126$ is capability
damage.} At $\alpha{=}84$ the harmful-specificity drop was selective. The primary evidence
is \emph{benign specificity} --- the direct mirror of the harmful actionability axis ---
which was preserved ($-0.005$, ns; \Cref{tab:regionsteer-selectivity}); consistent with
this, Group-C helpfulness was essentially unchanged (baseline $0.39\to0.37$;
\Cref{tab:regionsteer-guard}). At $\alpha{=}126$ the
harm dropped further ($-0.255$) but benign specificity fell with it ($-0.150$,
$p{=}8.5\!\times\!10^{-5}$) and Group-C utility collapsed (helpful $0.39\to0.15$, on-topic
$0.75\to0.10$, collapse $0.07\to0.21$); $\alpha{=}126$ is therefore the
\emph{capability-damage} regime, not a stronger result. The contrast with O is
instructive: O never degraded harmful specificity at any dose yet significantly reduced
\emph{benign} specificity from $\alpha{=}84$ on ($-0.107$, $-0.125$) --- the worst
possible profile, all cost and no safety benefit. Selective harm reduction was unique to
\texttt{tpl\_ctx} at $\alpha{=}84$; and even there the Group-C on-topic rate had begun to
soften ($0.75\to0.57$), placing $84$ at the \emph{edge} of the selective regime.
% RESOLVED (Group-C helpfulness @ alpha=84): print 0.37 (the selective operating point);
% 0.43 is the alpha=21/42 value. Selectivity rests on benign specificity (-0.005, ns) as
% the primary evidence; helpfulness (0.39 -> 0.37) is reported as consistent, not primary.

\paragraph{Precision: reduced actionability, not refusal.} At the selective point the
\emph{binary} inclusive-harm rate barely moved: T@$84$ delivered harm in 37/40 cases
(vs.\ 40/40 at baseline), and O/R/Rw stayed at 40/40 (R dipping to 38/40 only at
$\alpha{=}126$, within the straddler band of \Cref{ssec:determination}). T's inclusive
rate fell dose-dependently ($40\to39\to37\to30$ over $\alpha$) but never to refusal. The
honest reading is therefore that \texttt{tpl\_ctx} steering \emph{reduces the actionability
/ specificity} of harmful content without flipping it to a refusal; we do not claim
neutralisation.

% ----------------------------------------------------------------------------
\subsection{Synthesis}
\label{ssec:synthesis}

Behaviourally, the outcome of a DIJA attack was mostly determined by the scaffold ---
$\sim$97\% of draws at the paper's temperature went the way the scaffold pointed over an
unselected set, and the outcome was concordant with the $T{=}0$ generation in 14/15
cases --- yet the tokens carrying the harm were not written until step $\sim$98 of 128.
The $\sim$98 steps between prompt and commitment were therefore overwhelmingly
\emph{execution}, not deliberation: the trajectory mostly played out an outcome the
scaffold already fixed. This is a statement about the distribution of outcomes given a
scaffold, established by sampling; it is \emph{not} a claim that the model encodes the
decision in its step-0 representation, for which we have no evidence (the probe that
appeared to show it was reading arm identity, \Cref{ssec:placebo}). We say
\emph{determined by the scaffold}, not \emph{settled at prompt-encoding time in the
model's representations}.

\paragraph{One mechanism now accounts for every failed correction.} Detection looked
early and easy precisely \emph{because} the outcome was largely fixed by the input,
which is why AUC reached 1.000 and why that number means far less than it appears to.
Correction failed from the other side of the same fact, and the failures share a single
cause: every correction we tried intervenes on the \textbf{output region}, while the
determinant is the fixed scaffold text in the \textbf{prompt}. SAE feature zeroing and
residual-stream steering edit activations over the generated span; re-masking
(\Cref{ssec:remask}) regenerates the output tokens themselves --- the forward-process
lever a diffusion LM uniquely has. All three leave the harmful \emph{template} ---
immutable prompt tokens --- untouched, and it re-drives the same worksheet on every
refill, as arm~D's positional read shows directly. The detection--correction asymmetry
is thus \emph{explained, not merely observed}: the input nearly fixes the outcome, so
reading it is easy and moving it from inside the output region is structurally
impossible. Scope: within-template, within-model (LLaDA-8B-Instruct), the DIJA attack
family --- not a claim about correction in general.
% AUTHOR QUERY: cite keys for the SAE-feature-zeroing and residual-stream-steering
% experiments referenced here (the project's other two failed corrections)
% \citep{TODO-sae, TODO-steering}

\paragraph{The asymmetry has a causal complement --- and it partially rescues a
representational correction.} Result~4 established the asymmetry by elimination; the
region-steering experiment (\Cref{ssec:regionsteer}) closed it causally and directionally.
Steering the layer-16 representation of the \texttt{tpl\_ctx} scaffold region degraded
harmful specificity ($-0.123$ at $\alpha{=}84$, $-0.255$ at $126$), while the
\emph{identical} steer over the output region, a random output subset, or the wrapper did
nothing. Reading the input is easy and moving harm from \emph{inside the output region} is
structurally impossible --- but the scaffold's \emph{internal representation} is causally
and \emph{selectively} steerable at a bounded operating point ($\alpha{=}84$: harmful
actionability down, benign specificity and helpfulness preserved), with over-steering
($\alpha{=}126$) trading selectivity for capability damage. This constrains --- and may
partially rescue --- the paper's MOI steering framework: representation-level correction
was unavailable over the output region, but the scaffold representation is reachable at
\texttt{tpl\_ctx}, so such a correction is viable \emph{in principle} once retargeted
there --- subject to a selectivity/$\alpha$ ceiling, and only as reduced actionability,
never a refusal. Whether the \texttt{tpl\_ctx} result \emph{retargets} MOI or must instead
be \emph{reframed} against it turns on the exact token positions MOI's own steering wrote
to (see AUTHOR QUERY). It is an enabling mechanistic finding about \emph{where} the
correction handle lives, not a deployable safeguard.
% AUTHOR QUERY (PARKED --- awaiting co-author Nawaf): the expansion of ``MOI'' is
% UNRESOLVED, and its two candidates differ substantively: (a) this confirmation's synthesis
% originally glossed it as ``mechanism-of-injection correction''; (b) the existing paper's
% MOI/manifold section uses ``Manifold of Interest'' (the future-work steering framework).
% Do NOT commit either expansion. Pending from Nawaf: (1) the correct MOI expansion, and
% (2) which exact token positions the paper's existing MOI steering wrote to --- that
% determines whether the tpl_ctx result CONSTRAINS/RETARGETS MOI (if MOI steered
% output-region positions) or must be REFRAMED (if MOI already targeted prompt/scaffold
% positions).

\paragraph{The residual is the interesting part.} The $\sim$4\% of draws that go against
the scaffold (rising to $\sim$14\% at $T{=}0.7$, with degeneration and judge noise
excluded) are, by definition, the cases the scaffold did \emph{not} determine --- the
only population in which a mid-generation intervention has a forming choice to act on,
and the only one in which the representational question is even well-posed, since there
the prompt is byte-identical across outcomes. That population is concentrated in a
minority of scaffolds, not thinly spread.

\paragraph{Interpretation guard for any straddler probe.} Within a straddler the prompt
and the step-0 forward pass are byte-identical across outcome classes, which is what
makes the design leak-proof and forces $k_{\text{detect}} \ge 1$; but it also means the
two classes cannot differ for any reason present before sampling begins. The outcome
difference originates in sampling randomness alone.
\begin{quote}
A probe trained on straddler trajectories measures \textbf{when the trajectory becomes
committed to an outcome}, not when the model decides to comply. There is no decision in
the deliberative sense to locate: the ``decision'' is a coin flip that has already
happened in the sampler, and the probe watches its consequences propagate.
\end{quote}
This is sufficient for a runtime-monitoring claim (``by step $k$ the trajectory is
readably headed for harmful output'') but not for any mechanistic claim about
deliberation, and it does not generalise to the scaffolds that never flip. Write it as
\emph{commitment detection}, never as \emph{decision detection}.

% ----------------------------------------------------------------------------
\subsection{Limitations}
\label{ssec:limitations}

Decision-level $k_{\text{detect}}$ is not measurable in this regime, and that is a
property of the phenomenon: it requires same-prompt/different-outcome data, which at
$T{=}0$ does not exist (the outcome is deterministic) and at $T{=}0.2$ is rare (2.75\%)
and concentrated, so on the 25-scaffold set only 4 requests reach the $\ge 2$-minority
threshold a leave-one-group-out design needs. We do not report $\Delta$ as a number, and
we do not claim a pre-generation gate as a defense contribution (that is input filtering,
which needs no interpretability). Sub-token attribution of $k_{\text{commit}}$ is
unresolved in both directions. The label channel carries $\sim$2\% noise on borderline
texts, which caps an honest probe AUC at $\approx 0.96$ (a perfect probe against observed
labels at $\varepsilon = 0.02$), so an AUC of 1.000 should read as leakage. Result~4 is
10 requests, one dataset, a $T{=}0$ baseline regime, and a single (late) branch $k^\ast$;
the neutralisation result is descriptively unanimous but small, and the structural claim
--- the harmful template is fixed prompt text --- does not depend on $n$.
% AUTHOR QUERY: human-vs-judge agreement and Cohen's kappa are pending blinded
% adjudication (n=42, 2 annotators); insert the value once available and cross-ref
% \Cref{sec:judge-reconciliation}

\paragraph{Region-steering (\Cref{ssec:regionsteer}).} The effect is on
\emph{specificity / actionability}, not refusal: at the selective $\alpha{=}84$ the attack
still delivered inclusive harm in 37/40 cases, so this degrades how actionable the harmful
content is rather than neutralising it. The selective window is narrow --- $\alpha{=}84$
is selective but $\alpha{=}126$ is not (benign specificity $-0.150$, Group-C utility
collapse), and even at $84$ the Group-C on-topic rate had begun to soften
($0.75\to0.57$) --- so the operating point is load-bearing and would need per-model,
per-attack validation before any deployment claim. The result is $n{=}40$, one model
(LLaDA-8B-Instruct), one attack family (DIJA), and a single steering axis (the layer-16
injection direction). We claim selective prompt-region \emph{steerability} --- an enabling
mechanism --- not a working defense.
